\documentclass[10pt,a4paper]{amsart}
\usepackage[margin=19mm]{geometry}
\usepackage{amsmath,amssymb,mathtools}
\usepackage[scr=rsfs,cal=euler]{mathalfa}
\usepackage{xfrac}
\usepackage[unicode,hidelinks,bookmarks=false]{hyperref}
\newcommand{\ie}{i.e.\ }
\newcommand{\cf}{cf.\ }
\newcommand{\resp}{respectively\ }
\newcommand{\Subsection}{\subsection}
\providecommand{\nobreakdash}{\nobreak\hbox{-}}
\setlength{\emergencystretch}{2em}
\multlinegap0pt
\usepackage{amssymb}
\newtheorem{defn}{Definition}
\newtheorem{thm}{Theorem}
\newcommand{\R}{\mathbb R}
\newcommand{\ds}{\displaystyle}
\title{Canonical parameters on marginally trapped surfaces in the Minkowski 4-space}
\author{Miroslav Maksimovi\'c, Velichka Milousheva}
\date{Source: arXiv:2605.08367v1 (CC BY 4.0)}
\begin{document}
\maketitle

\noindent\emph{Source and licence.} Canonical parameters on marginally trapped surfaces in the Minkowski 4-space. Version arXiv:2605.08367v1 (CC BY 4.0), distributed by arXiv under the Creative Commons Attribution 4.0 International licence, http://creativecommons.org/licenses/by/4.0/, as recorded in the arXiv licence metadata for this exact version and shown on the abstract page https://arxiv.org/abs/2605.08367v1. Full licence notice: \texttt{licenses/arxiv-2605.08367-CC-BY-4.0.txt}.\par\smallskip

\noindent\textit{Editorial scope.} Counted unit: the article's thm eq:ExistenceCanonicalParameters (source0.tex lines 319--321). The article's complete original proof of it is delivered with it (source0.tex lines 323--403, one \texttt{proof} environment of 81 lines). No mathematics is written, repaired or re-derived: the statement and the proof are the article's own lines, in the article's own notation, and no retained line is substituted or re-flowed.  Retained uncounted as the source-local context the proof needs: lines 295--304; lines 311--313.  Nothing was omitted from any retained span, so the package carries no omission record.  No retained line contains a citation, so the wrapper carries no bibliography and no bibliography entry.  No retained line contains a figure environment, an image-inclusion command or drawing code, so no image is delivered and the package declares no asset dependency.  Editorial formatting only, in addition: an AMS \texttt{amsart} preamble that restates the article's own declarations and macros used by the retained lines, namely the environments \texttt{defn}, \texttt{thm} on separate counters, together with the standard packages those lines need.  The retained environments print in this document as Definition 1, Theorem 1 in order of appearance, whereas the article prints the same items with the numbering of its own full document.  The complete native source of the article as distributed in the arXiv e-print for this exact version is bundled unchanged as \texttt{sources/200149/source0.tex}.
\begin{equation} \label{eq:phi-psi}
	\begin{array}{l} 
		\vspace{2mm}
		\varphi(u) = \ds\sqrt Ee^{-\ds\int_{v_0}^v\left(\phi_1+\phi_2\frac{\sqrt G}{\sqrt E}\right)dv
			-\int_{u_0}^u\left(\phi_3\frac{\sqrt E}{\sqrt G}+\phi_4\right)(u,v_0)du+c_1}\ ;  \\
		\vspace{2mm}
		\psi(v) = \ds\sqrt Ge^{-\ds\int_{u_0}^u\left(\phi_3\frac{\sqrt E}{\sqrt G}+\phi_4\right)du
			-\int_{v_0}^v\left(\phi_1+\phi_2\frac{\sqrt G}{\sqrt E}\right)(u_0,v)dv+c_2} \ .
	\end{array}   
\end{equation}

\begin{defn} \label{D:def-can}
Let $M^2$ be a marginally trapped surface in $\R^4_1$ parametrized by principal parameters $(u,v)$ and  $K^2+\varkappa^2 \neq \mu^2$. If the functions $\varphi(u)$ and $\psi(v)$ defined by \eqref{eq:phi-psi} are equal to 1, then we say that the parameters $(u,v)$ are {\it canonical principal parameters} of the surface.
\end{defn}

\begin{thm} \label{eq:ExistenceCanonicalParameters}
	Each marginally trapped surface with  $K^2+\varkappa^2 \neq \mu^2$ locally admits canonical principal parameters.
\end{thm}

\begin{proof}
	For a given  pair of principal parameters $(u,v)$ and an arbitrary point $(u_0,v_0)$, we introduce new parameters $\bar u$, $\bar v$ defined by
	\begin{equation*} \label{eq:defCanonParam}
		\begin{array}{l}
			\bar u=\ds\int_{u_0}^u\sqrt Ee^{-\ds\int_{v_0}^v\left(\phi_1+\phi_2\frac{\sqrt G}{\sqrt E}\right)dv
				-\int_{u_0}^u\left(\phi_3\frac{\sqrt E}{\sqrt G}+\phi_4\right)(u,v_0)du+c_1} +u_0 \vspace{2mm} \ ;\\
			\bar v=\ds\int_{v_0}^v\sqrt Ge^{-\ds\int_{u_0}^u\left(\phi_3\frac{\sqrt E}{\sqrt G}+\phi_4\right)du
				-\int_{v_0}^v\left(\phi_1+\phi_2\frac{\sqrt G}{\sqrt E}\right)(u_0,v)dv+c_2} +v_0 \ ,
		\end{array}
	\end{equation*}  
	or equivalently,
	\begin{equation} \label{eq:defCanonParam1}
	\bar u=\int_{u_0}^u \varphi(u) du +u_0 \ ;   \qquad  \bar v=\int_{v_0}^v \psi(v) dv +v_0 ,
		\end{equation}  
	where the functions $\varphi(u)$ and $\psi(v)$ are given by \eqref{eq:phi-psi}. Therefore,  $\bar u=\bar u(u)$ and $\bar v=\bar v(v)$ hold, which means that the new  parameters 
	$(\bar u, \bar v)$ are also principal. Obviously, $\bar u_0=\bar u(u_0)=u_0$, 
	$\bar v_0=\bar v(v_0)=v_0$. Differentiating  equations \eqref{eq:defCanonParam1} gives
	\begin{equation*} 
		\begin{array}{l} 
			\vspace{2mm}
	\bar u_u = \varphi(u) =\sqrt Ee^{-\ds\int_{v_0}^v\left(\phi_1+\phi_2\frac{\sqrt G}{\sqrt E}\right)dv
		-\int_{u_0}^u\left(\phi_3\frac{\sqrt {E}}{\sqrt {G}}+\phi_4\right)(u,v_0)du+c_1} \ ;\\
			\vspace{2mm}
	\bar u_v = 0 \ ;\\
	\vspace{2mm}
	\bar v_u = 0 \ ;\\
	\vspace{2mm}
	\bar v_v = \psi(v) =\sqrt Ge^{-\ds\int_{u_0}^u\left(\phi_3\frac{\sqrt E}{\sqrt G}+\phi_4\right)du
		-\int_{v_0}^v\left(\phi_1+\phi_2\frac{\sqrt G}{\sqrt E}\right)(u_0,v)dv+c_2} \ .
	\end{array}   
	\end{equation*}
	
	
	It can easily be checked that  under the change of the parameters $\bar u=\bar u(u)$, $\bar v= \bar v(v) $ we have the relations:
	$$\sqrt {E} = \sqrt {\bar E} \, {\bar u}_u; \quad \sqrt {G} = \sqrt {\bar G} \, {\bar v}_v,$$
	from where it follows that
	\begin{equation*} 
		\begin{array}{ll} 
			\vspace{2mm}
			\phi_1 (u,v) = {\bar \phi_1} (\bar u(u),\bar v(v)) \, {\bar v}_v;  & \quad \phi_2 (u,v) = {\bar \phi_2} (\bar u(u),\bar v(v)) \, {\bar u}_u;  \\
			\vspace{2mm}
			\phi_3 (u,v) = {\bar \phi_3} (\bar u(u),\bar v(v)) \, {\bar v}_v;  & \quad \phi_4 (u,v) = {\bar \phi_4} (\bar u(u),\bar v(v)) \, {\bar u}_u.
		\end{array}   
	\end{equation*}
The last equations imply 
	$$
	\ds (\phi_1+\phi_2\frac{\sqrt {G}}{\sqrt {E}})(u,v)=(\bar \phi_1+\bar \phi_2\frac{\sqrt {\bar G}}{\sqrt {\bar E}})(\bar u(u),\bar v(v))\frac{d\bar v}{dv} \ ;
	$$
	$$
	\ds (\phi_3\frac{\sqrt {E}}{\sqrt {G}}+\phi_4)(u,v)=(\bar \phi_3\frac{\sqrt {\bar E}}{\sqrt {\bar G}}+\bar \phi_4)(\bar u(u),\bar v(v))\frac{d\bar u}{du} \ .
	$$
By calculations we further obtain
	\begin{equation}\label{eq:3.6}
		\begin{split}
			& e^{\ds\int_{v_0}^v\left(\phi_1+\phi_2\frac{\sqrt G}{\sqrt E}\right)(u,v)dv
				+\int_{u_0}^u\left(\phi_3\frac{\sqrt {E}}{\sqrt {G}}+\phi_4\right)(u,v_0)du} = \\
			& =e^{\ds\int_{v_0}^v\left(\bar \phi_1+\bar \phi_2\frac{\sqrt {\bar G}}{\sqrt {\bar E}}\right)(\bar u(u),\bar v(v))\frac{d\bar v}{dv}{dv}
				+\int_{u_0}^u\left(\bar \phi_3\frac{\sqrt {\bar E}}{\sqrt {\bar G}}+\bar \phi_4\right)(\bar u(u),\bar v(v_0))\frac{d\bar u}{du}du} = \\
			& =e^{\ds\int_{\bar v_0}^{\bar v}\left(\bar \phi_1+\bar \phi_2\frac{\sqrt {\bar G}}{\sqrt {\bar E}}\right)(\bar u,\bar v)d\bar v
				+\int_{\bar u_0}^{\bar u}\left(\bar \phi_3\frac{\sqrt {\bar E}}{\sqrt {\bar G}}+\bar \phi_4\right)(\bar u,\bar v_0)d\bar u} \ .
		\end{split}
	\end{equation}
		Using the previous results we have
	$$
	\sqrt{\bar E}=\frac{\sqrt{E}}{\bar u_u}=e^{\ds\int_{v_0}^v\left(\phi_1+\phi_2\frac{\sqrt G}{\sqrt E}\right)dv
		+\int_{u_0}^u\left(\phi_3\frac{\sqrt {E}}{\sqrt {G}}+\phi_4\right)(u,v_0)du-c_1}
	$$
	
	$$
	=e^{\ds\int_{\bar v_0}^{\bar v}\left(\bar \phi_1+\bar \phi_2\frac{\sqrt {\bar G}}{\sqrt {\bar E}}\right)(\bar u,\bar v)d\bar v
		+\int_{\bar u_0}^{\bar u}\left(\bar \phi_3\frac{\sqrt {\bar E}}{\sqrt {\bar G}}+ \bar \phi_4\right)(\bar u,\bar v_0)d\bar u-c_1}
	$$
	and hence $\bar\varphi (\bar u)=1$, where the function $\bar\varphi$ is given by
$$
\ds\bar\varphi (\bar u) = \sqrt{ \bar E} e^{\ds-\int_{\bar{v}_0}^{\bar v}\left(\bar\phi_1+\bar\phi_2\frac{\sqrt{\bar G}}{\sqrt{\bar E}}\right)d\bar v
	-\int_{\bar{u}_0}^{\bar u}\left(\bar\phi_3\frac{\sqrt{\bar E}}{\sqrt{\bar G}}+\bar\phi_4\right)(\bar u,\bar{v}_0)d\bar u+c_1}.
$$

In a  similar way, we also obtain that $\bar\psi(\bar v)=1$, where the function $\bar\psi$ is defined analogously to function $\psi$, but with respect to the parameters $(\bar u,\bar v)$. Thus, according to Definition \ref{D:def-can}, the parameters $(\bar u,\bar v)$ are canonical.
	
\end{proof}

\end{document}
